\chapter{Key Methods and Core Mechanisms}


\section{Controlled Execution and Evaluation Pipeline}

The key design is not simply running code, but decomposing each submission into a reviewable, traceable, and explainable process. After submission, the backend validates identity and assignment legality, stores submission records, and then performs execution, analysis, grading, and feedback generation in a fixed order to avoid unexplained result fluctuations over time.

At execution level, user code is run in a controlled environment with constraints on runtime duration and resource usage, while outputs are collected systematically. This provides two direct teaching benefits: instructors do not need to maintain multiple local runtime environments, and students receive more consistent execution outcomes, reducing disputes such as "works locally but fails on the platform."

Execution outcomes are stored in structured form, including stdout, stderr, and exit status. These records support both immediate grading and later instructor review, meaning each score is backed by process-level evidence rather than a single final number.

\section{Code Quality Analysis and Fairness Control}

Using hidden-test pass rate alone makes it difficult to distinguish between two cases: functionally correct code with poor quality, and conceptually reasonable code with implementation details not yet complete. To address this, the system adopts a dual-track mechanism: final score for functional correctness, and static-analysis score for coding quality, structure, and process quality.

The static-analysis module scans code under rule-based checks, extracts potential issues, and applies corresponding deductions. This score supplements the final score instead of replacing it. As a result, students can understand not only whether a solution is correct, but also where it can be improved; instructors can also identify class-wide patterns more efficiently.

To reduce scoring distortion, the system includes template/blank-submission detection. If submitted code is highly similar to the starter template or lacks effective newly added logic, blank-submission detection is triggered, static-analysis scores are reduced, and warnings are displayed. This mechanism helps prevent inflated scores from superficial submissions and makes grading more aligned with actual student effort.

\section{RAG-Based Feedback and Teaching Alignment}

In the feedback stage, student code is not sent directly to a general-purpose model. Instead, relevant course materials are retrieved first, and retrieval outputs are combined with submission context for generation. The purpose is to keep feedback centered on lecture content, course terminology, and assignment objectives, reducing generic or out-of-scope advice.

The process can be summarized as "retrieve first, generate second." The system retrieves relevant teaching fragments based on assignment topic and code features, then combines runtime results, static-analysis conclusions, and grading information to produce structured feedback. Feedback typically includes issue localization, revision suggestions, and actionable next steps.

This design is consistent with findings in prior work: RAG improves factual grounding by making retrieved evidence part of the generation condition, not just an auxiliary reference \cite{ref3_lewis2020}. For code-related outputs, this thesis also prioritizes executable signals, including test outcomes and runtime traces, over pure text overlap, matching evidence from Codex and large-scale synthesis studies that functional correctness is a more reliable target than n-gram similarity metrics such as BLEU \cite{ref5_chen2021,ref6_austin2021}.

To ensure stability in classroom use, a fallback path is preserved: when retrieval is insufficient or model services fluctuate, the system can still return basic rule-based feedback so that page readability and workflow continuity are maintained. In short, RAG improves feedback quality without becoming a single point of failure in the main pipeline.

